\section{ウォリスの公式}\label{節:ウォリス公式}

前節の方法を三角関数のべきに使ってみよう.
\[
 I_n=\int_0^{\frac{\pi}{2}}\sin^n x\dd x
 \qquad(n=0,1,2,\ldots)
\]
を考える.最初の2つは
\[
 I_0=\frac{\pi}{2},\qquad I_1=1
\]
である.\,$n\geqq2$の$I_n$を,指数の小さい積分に帰着させることを目標とする.

$\sin^n x$を$\sin^{n-1}x\cdot\sin x$と分け,後ろの$\sin x$を積分する側にして部分積分する.微分された$\sin^{n-1}x$から$\cos x$が現れ,積分の中に$\cos^2x$ができる.これを$1-\sin^2x$に書き直せば,再び$\sin$のべきだけの積分になる.

ここで1つ断りをいれておこう.\,$n\geqq2$なので$\sin^{n-1}0=0$であり,また$\cos\frac\pi2=0$である.したがって部分積分の端点の項は両端とも$0$になる.
\begin{align*}
 I_n
 &=\left[-\sin^{n-1}x\cos x\right]_0^{\frac{\pi}{2}}
 +(n-1)\int_0^{\frac{\pi}{2}}\sin^{n-2}x\cos^2x\dd x\\
 &=(n-1)\int_0^{\frac{\pi}{2}}\sin^{n-2}x(1-\sin^2x)\dd x\\
 &=(n-1)(I_{n-2}-I_n)
\end{align*}
最後の式の$I_n$を左辺に移項して$n$で割ると
\begin{equation}\label{式:ウォリス漸化式}
 I_n=\frac{n-1}{n}I_{n-2}\qquad(n\geqq2)
\end{equation}
を得る.

この漸化式は,\,$I_n$を2つ前の$I_{n-2}$と結ぶ.指数が$2$ずつ下がるので,偶数番目は$I_0$へ,奇数番目は$I_1$へつながる.偶数番目だけを取り出して$u_m=I_{2m}$と置くと,\,$u_m=\frac{2m-1}{2m}u_{m-1}$となり,隣り合う項の比が$m$の式で与えられる階比型になる.よって比を順に掛けてやれば
\begin{align*}
 I_{2m}
 &=\frac{2m-1}{2m}I_{2m-2}\\
 &=\frac{2m-1}{2m}\cdot\frac{2m-3}{2m-2}I_{2m-4}\\
 &\qquad\vdots\\
 &=\frac{2m-1}{2m}\cdot\frac{2m-3}{2m-2}\cdots\frac12I_0
\end{align*}
となる.奇数番目も同じく$I_1$まで下げればよい.

\begin{theorem}{ウォリスの積分公式}
$m\geqq1$に対して
\begin{equation}\label{式:ウォリス偶数}
 I_{2m}
 =\frac{1\cdot3\cdot5\cdots(2m-1)}{2\cdot4\cdot6\cdots2m}
   \cdot\frac{\pi}{2},
\end{equation}
\begin{equation}\label{式:ウォリス奇数}
 I_{2m+1}
 =\frac{2\cdot4\cdot6\cdots2m}{3\cdot5\cdot7\cdots(2m+1)}.
\end{equation}
\end{theorem}

偶数の積と奇数の積を短く書くために
\[
 (2m)!!=2\cdot4\cdots2m,\qquad
 (2m-1)!!=1\cdot3\cdots(2m-1)
\]
と定義することがある.これは\textbf{二重階乗}と呼ばれ,階乗を2回行うという意味ではない.この記号なら
\[
 I_{2m}=\frac{(2m-1)!!}{(2m)!!}\cdot\frac{\pi}{2},\qquad
 I_{2m+1}=\frac{(2m)!!}{(2m+1)!!}
\]
となる.

\begin{bookpoint}
\textbf{POINT}\par
$\sin^nx=\sin^{n-1}x\cdot\sin x$と分けて部分積分し,\,$\cos^2x=1-\sin^2x$で$I_n$を戻して移項する.\,$I_n=\frac{n-1}{n}I_{n-2}$を$I_0$か$I_1$まで繰り返す.
\end{bookpoint}

\begin{example}{ウォリス積分}
次の2つの積分を求めよ.
\[
 \int_0^{\frac{\pi}{2}}\sin^6x\dd x,\qquad
 \int_0^{\frac{\pi}{2}}\sin^7x\dd x.
\]
\end{example}

\paragraphbf{方針}
漸化式\eqref{式:ウォリス漸化式}で,偶数は$I_0=\frac\pi2$まで,奇数は$I_1=1$まで指数を下げる.

\paragraphbf{解答}
$\int_0^{\frac{\pi}{2}}\sin^6x\dd x=\frac{5\pi}{32}$,\qquad
$\int_0^{\frac{\pi}{2}}\sin^7x\dd x=\frac{16}{35}$

\paragraphbf{解法}
漸化式\eqref{式:ウォリス漸化式}を繰り返すと
\begin{align*}
 I_6&=\frac56\cdot\frac34\cdot\frac12\cdot\frac{\pi}{2}
     =\frac{5\pi}{32},\\
 I_7&=\frac67\cdot\frac45\cdot\frac23\cdot1
     =\frac{16}{35}
\end{align*}
となる.偶数と奇数で最後に戻る値が違うことが,答えに$\pi$が現れるかどうかの違いになっている.

積分区間が$\left[0,\frac{\pi}{2}\right]$であることも大切である.任意の区間では,部分積分の端点の項が同じように消えるとは限らない.

\subsection{積分の比}

公式は積分の値を積で表したものである.今度は逆に,積分の大小を使ってその積の極限を求める.まず$I_n$の値を並べてみよう.
\begin{center}
\begin{tabular}{cccccccccc}
\toprule
$n$ & $0$ & $1$ & $2$ & $3$ & $4$ & $5$ & $6$ & $7$ & $8$\\
\midrule
$I_n$ & $1.5708$ & $1$ & $0.7854$ & $0.6667$ & $0.5890$ & $0.5333$ & $0.4909$ & $0.4571$ & $0.4295$\\
\bottomrule
\end{tabular}
\end{center}
$I_n$は減少し,しかも隣り合う値の差が次第に小さくなっていく.隣どうしの比$\frac{I_{n+1}}{I_n}$は$1$へ近づくだろう.

実際,\,$0\le\sin x\le1$だから,指数を大きくすると被積分関数は小さくなる.したがって
\[
 I_{n+2}\le I_{n+1}\le I_n
\]
である.各辺を$I_n>0$で割り,式\eqref{式:ウォリス漸化式}を使うと
\[
 \frac{n+1}{n+2}
 =\frac{I_{n+2}}{I_n}
 \le\frac{I_{n+1}}{I_n}\le1
\]
となる.両端が$1$へ近づくため
\begin{equation}\label{式:ウォリス隣接比}
 \lim_{n\to\infty}\frac{I_{n+1}}{I_n}=1
\end{equation}
である.各$I_n$の極限を先に求めなくても,比の極限は分かるのだ.

ここで,次の有限の積を作る.
\[
 P_m=\prod_{k=1}^{m}\frac{(2k)^2}{(2k-1)(2k+1)}
\]
$\prod$は,指定された範囲の項をすべて掛ける記号である.最初の数項を書けば
\[
 P_m=\frac{2\cdot2}{1\cdot3}\cdot
      \frac{4\cdot4}{3\cdot5}\cdots
      \frac{2m\cdot2m}{(2m-1)(2m+1)}
\]
となる.分子の偶数の積は$\{(2m)!!\}^2$,分母の奇数の積は$(2m-1)!!\,(2m+1)!!$である.一方,式\eqref{式:ウォリス偶数}と式\eqref{式:ウォリス奇数}の比を取ると
\begin{align*}
 \frac{I_{2m+1}}{I_{2m}}
 &=\frac{2}{\pi}\cdot
   \frac{\{(2m)!!\}^2}{(2m-1)!!\,(2m+1)!!}\\
 &=\frac{2}{\pi}P_m
\end{align*}
となる.左辺は$1$へ近づくので,次の極限が得られる.
\begin{equation}\label{式:ウォリスの積}
 \boxed{
 \lim_{m\to\infty}\prod_{k=1}^{m}
 \frac{4k^2}{4k^2-1}=\frac{\pi}{2}}
\end{equation}
これが\textbf{ウォリスの積}である.無限個の数を一度に掛けたと考えるのではなく,有限の積$P_m$の数列の極限として定義している.

\paragraphbf{円周率を挟む}
証明の途中の不等式を残しておけば,近似の誤差も評価できる.\,$n=2m$とすると
\[
 \frac{2m+1}{2m+2}
 \le\frac{I_{2m+1}}{I_{2m}}\le1
\]
だから
\[
 \frac{2m+1}{2m+2}\cdot\frac{\pi}{2}
 \le P_m\le\frac{\pi}{2}
\]
である.各辺に正の数を掛けて$\pi$について解くと
\[
 2P_m\le\pi\le\frac{2m+2}{2m+1}\cdot2P_m
\]
となる.実際に計算してみよう.
\begin{center}
\begin{tabular}{ccc}
\toprule
$m$ & $2P_m$ & $\frac{2m+2}{2m+1}\cdot2P_m$\\
\midrule
$1$ & $2.6667$ & $3.5556$\\
$10$ & $3.0677$ & $3.2138$\\
$100$ & $3.1338$ & $3.1494$\\
$1000$ & $3.1408$ & $3.1424$\\
\bottomrule
\end{tabular}
\end{center}
$\pi=3.14159\ldots$は,どの$m$でも左右の値の間にある.有限の計算の結果が,どれほど確かな範囲を与えるかまで分かったのだ.ただし,区間の幅は$\frac{2P_m}{2m+1}$であり,\,$\frac{\pi}{2m+1}$より小さいものの,\,$m$に反比例する程度でしか縮まらない.\,$m=1000$でも,確定するのは$3.14$までである.
%
\subsection{二項係数}

偶数の公式の分数は,階乗で書き直せる.偶数の積からは$2$を1つずつくくり出せるので
\[
 (2m)!!=2^m m!
\]
である.また$(2m)!$は偶数の積と奇数の積に分かれるので
\[
 (2m-1)!!=\frac{(2m)!}{(2m)!!}=\frac{(2m)!}{2^m m!}
\]
となる.これらを式\eqref{式:ウォリス偶数}に代入すると
\begin{equation}\label{式:ウォリスと二項係数}
 I_{2m}=\frac{\pi}{2}\cdot\frac{(2m)!}{2^{2m}(m!)^2}
 =\frac{\pi}{2}\cdot\frac{1}{4^m}\binom{2m}{m}
\end{equation}
を得る.ここで$\binom{2m}{m}=\frac{(2m)!}{(m!)^2}$は二項係数であり,\,$2m$個から$m$個を選ぶ組合せの数${}_{2m}\mathrm{C}_m$に等しい.

二項係数$\binom{2m}{m}$は$m$とともに急速に大きくなる.では,\,$4^m$で割った値はどうなるだろうか.\,$\sqrt m$を掛けて調べてみよう.
\begin{center}
\begin{tabular}{cccccc}
\toprule
$m$ & $1$ & $10$ & $100$ & $1000$ & $10000$\\
\midrule
$\frac{\sqrt m}{4^m}\binom{2m}{m}$ & $0.5$ & $0.5572$ & $0.5635$ & $0.5641$ & $0.5642$\\
\bottomrule
\end{tabular}
\end{center}
値は$0.564$付近へ近づいているように見える.この値を,ウォリスの積分から求めよう.

偶数の公式と奇数の公式を掛けると,分子と分母の偶数がすべて約分されて
\[
 I_{2m}I_{2m+1}
 =\frac{1\cdot3\cdots(2m-1)}{3\cdot5\cdots(2m+1)}\cdot\frac{\pi}{2}
 =\frac{\pi}{2(2m+1)}
\]
となる.これと式\eqref{式:ウォリス隣接比}の$\frac{I_{2m+1}}{I_{2m}}\to1$を組み合わせれば
\begin{align*}
 mI_{2m}^2
 &=m\,I_{2m}I_{2m+1}\cdot\frac{I_{2m}}{I_{2m+1}}\\
 &=\frac{m\pi}{2(2m+1)}\cdot\frac{I_{2m}}{I_{2m+1}}
 \longrightarrow\frac{\pi}{4}
\end{align*}
となる.\,$I_{2m}>0$なので正の平方根を取ることができ,
\[
 \sqrt m\,I_{2m}\longrightarrow\frac{\sqrt\pi}{2}
\]
である.式\eqref{式:ウォリスと二項係数}より$\frac{\sqrt m}{4^m}\binom{2m}{m}=\frac2\pi\sqrt m\,I_{2m}$なので
\begin{equation}\label{式:中央二項係数の極限}
 \lim_{m\to\infty}\frac{\sqrt m}{4^m}\binom{2m}{m}
 =\frac1{\sqrt\pi}
\end{equation}
を得る.表の$0.564$付近の値は,\,$\frac1{\sqrt\pi}=0.56418\ldots$だったのだ.階乗を含む数列の大きさが,三角関数の積分を経由して分かったことになる.

\paragraphbf{硬貨を投げる確率}
硬貨を$2m$回投げて表がちょうど$m$回出る確率は,\,$\frac1{4^m}\binom{2m}{m}$である.式\eqref{式:中央二項係数の極限}によれば,この確率は$m$が大きいとき$\frac{1}{\sqrt{\pi m}}$に近い.例えば$200$回投げて表がちょうど$100$回出る確率は,\,$\frac1{\sqrt{100\pi}}=0.0564\ldots$に近く,実際の値は$0.0563\ldots$である.表と裏が半分ずつになることが最も起こりやすいとはいえ,ちょうど半分になる確率そのものは$0$へ近づいていく.
%
